% ═══════════════════════════════════════════════════════════════════════
%  Macros mathématiques pour la version web (MathJax)
%
%  Équivalents MathJax des macros de commandesJS.sty et de macros-livre.tex.
%  Le PDF n'utilise PAS ce fichier : il lit directement les originaux.
%  Toute nouvelle macro utilisée dans une formule doit donc être ajoutée
%  à la fois dans macros-livre.tex (PDF) et ici (web).
%  \num, \SI et \ang (siunitx) sont traités par le filtre, pas par MathJax.
% ═══════════════════════════════════════════════════════════════════════

% ── Ensembles de nombres (commandesJS) ─────────────────────────────────
\newcommand{\R}{\mathbb{R}}
\newcommand{\C}{\mathbb{C}}
\newcommand{\Z}{\mathbb{Z}}
\newcommand{\Q}{\mathbb{Q}}
\newcommand{\N}{\mathbb{N}}
\newcommand{\Ns}{\mathbb{N}^{\star}}
\newcommand{\Rs}{\mathbb{R}^{\star}}
\newcommand{\Rp}{\left[0, \infty\right[}
\newcommand{\Rps}{\left]0, \infty\right[}

% ── Constantes, texte en exposant ──────────────────────────────────────
\newcommand{\eu}{\mathrm{e}}
\newcommand{\iu}{\mathrm{i}\mkern1mu}
\newcommand{\uppi}{\mathrm{\pi}}
\newcommand{\upgamma}{\mathrm{\gamma}}
\newcommand{\ap}[1]{^{\mathrm{#1}}}
\newcommand{\ped}[1]{_{\mathrm{#1}}}

% ── Intervalles (paquet interval : crochets « à la française ») ────────
\newcommand{\intff}[2]{\mathopen{[}#1, #2\mathclose{]}}
\newcommand{\intoo}[2]{\mathopen{]}#1, #2\mathclose{[}}
\newcommand{\intof}[2]{\mathopen{]}#1, #2\mathclose{]}}
\newcommand{\intfo}[2]{\mathopen{[}#1, #2\mathclose{[}}

% ── Analyse, probabilités, ensembles ───────────────────────────────────
\newcommand{\evint}[3]{\left[#1\right]_{#2}^{#3}}
\newcommand{\ol}[1]{\overline{#1}}
\newcommand{\zbar}{\overline{z}}
\newcommand{\RE}{\operatorname{Re}}
\newcommand{\IM}{\operatorname{Im}}
\newcommand{\Arg}{\operatorname{Arg}}
\newcommand{\P}[1]{\mathbb{P}\left(#1\right)}
\newcommand{\Es}[1]{\mathbb{E}\left(#1\right)}
\newcommand{\Var}[1]{\operatorname{var}\left(#1\right)}
\newcommand{\Cov}[1]{\operatorname{cov}\left(#1\right)}
\newcommand{\card}[1]{\operatorname{Card}\left(#1\right)}
\newcommand{\comp}[1]{#1^{\mathsf{c}}}
\newcommand{\pgcd}{\operatorname{pgcd}}
\newcommand{\ppcm}{\operatorname{ppcm}}
\newcommand{\proj}{\operatorname{proj}}
\newcommand{\egdef}{\overset{\scriptsize\text{déf}}{=}}
\newcommand{\Dom}[1]{\mathcal{D}_{#1}}
\newcommand{\Img}[1]{\mathcal{I}_{#1}}
\newcommand{\survie}[1]{\overline{\mkern-2mu#1\mkern2mu}}

% ── Algèbre linéaire (commandesJS) ─────────────────────────────────────
\newcommand{\matdeux}[4]{\begin{pmatrix}#1&#2\\#3&#4\end{pmatrix}}
\newcommand{\vdeux}[2]{\begin{pmatrix}#1\\#2\end{pmatrix}}
\newcommand{\vtrois}[3]{\begin{pmatrix}#1\\#2\\#3\end{pmatrix}}

% ── Raccourcis du volume (macros-livre.tex) ──────────────────────────
\newcommand{\boitevide}{\square}
\newcommand{\Cs}{\mathbb{C}^{\star}}
\newcommand{\coloneqq}{\mathrel{:=}}
